\section*{v1.2.0 Consolidation Core}
\addcontentsline{toc}{section}{v1.2.0 Consolidation Core}

\paragraph{Normalization note.} This block is retained as a \emph{historical consolidation layer}. In the normalized grouped architecture of Version~\docversion, its active generative content is absorbed into the normative layer $\mathfrak G$; the present section therefore records the earlier consolidation surface and should no longer be read as an independent foundational core competing with the later normalized structural kernel.


This version closes the gap between the previously inserted roadmap and the
already existing structural sections of the monolith. Its purpose is to state,
in one contiguous place, the working discipline of the next release, the core
closure chain, the probe-sector architecture, the first-layer master-equation
logic, and the normalized tool catalogue. The guiding rule remains unchanged:
\emph{no shortening, no silent merging, no replacement of the leading line;
only sharpening, explicit classification, and append-only completion.}

\subsection*{1. Structural reading discipline}
Every major claim used below is assigned one of four statuses.
\begin{enumerate}[leftmargin=2em,label=(\alph*)]
  \item \textbf{Structurally forced:} the statement follows from the ground
  condition, the triolic organization, orthogonality, and closure.
  \item \textbf{Projectively read:} the statement is not primary, but appears
  after passage to a chosen observational or effective sector.
  \item \textbf{Heuristically interpreted:} the structural content is stable,
  while its physical naming is still a disciplined reading rather than a final
  derivation.
  \item \textbf{Open:} the location of the claim is fixed, but its expanded
  proof or explicit computation is postponed to the next lemma layer.
\end{enumerate}

This fourfold discipline is essential because the framework is intended to be a
\emph{structure-first} proof architecture. In particular, equations are not to
be introduced as autonomous starting axioms when they can instead be read as
forced compression formulas of a deeper closure mechanism.

\subsection*{2. Ground condition, HC, and closure}
The working core of v1.2.0 is the explicit chain
\[
\text{Ground Condition} \Longrightarrow \text{HC} \Longrightarrow \text{Closure}.
\]
Here the ground condition is the requirement that the underlying triolic body is
not free to unfold into arbitrary independent sectors. Rather, the threefold
organization is constrained so that one component remains orthogonal to the
other two, the admissible units are treated in the standing imaginary
convention, and every physical readout is already phase-shifted with respect to
the underlying structural frame.

Under this condition, the Heisenberg compensator is not an external correction
patched in after the fact. It is the minimal balancing mechanism that becomes
necessary once a common body must be represented consistently in dual or
phase-shifted ways without allowing the structure to drift apart. In the logic
of the present paper, HC therefore carries three simultaneous tasks: it
re-identifies dual or shifted representatives, it suppresses unstable escape
modes, and it transports the system toward a light/collapse-compatible output.

Closure is then the global consequence of this balancing action. Once the
compensator has become active, admissible dynamics are no longer free to leave
the confined regime; instead they remain trapped within the higher coupled
architecture represented in the paper by the septinion/seption perspective.
This is the sense in which HC is emergent, while closure is enforced.

\paragraph{Proposition A (ground condition induces HC).}
Given the triolic structural core, the orthogonality constraint, and the
phase-shifted measurement principle, a compensating identification mechanism is
forced whenever the same body is to admit more than one consistent appearance.
That mechanism is the Heisenberg compensator.

\paragraph{Proof sketch.}
Without compensation, dual or shifted presentations of the same body would
produce uncontrolled representation drift. Because the framework demands one
common object across all sectors, a balancing map is required. This balancing
map is precisely the HC in structural form.

\paragraph{Proposition B (HC induces closure).}
Whenever HC acts as the balancing constraint of the triolic system, the
admissible dynamics are confined to a closure-stable regime. In particular,
sectors may transform, couple, or be read out differently, but they cannot
escape the common structural container without violating the ground condition.

\paragraph{Proof sketch.}
HC identifies and neutralizes precisely those deviations that would split the
common body into incompatible local descriptions. The remaining degrees of
freedom therefore assemble into a closed class. This class is the closure
regime read throughout the monolith.

\paragraph{Protected invariants.}
At the level of v1.2.0, the most important protected invariants are:
(i) the persistence of one common underlying body, (ii) orthogonality of the
third triolic direction relative to the first two, (iii) frame-shift discipline
between structure and observer description, and (iv) confinement of admissible
dynamics to the closure-compatible sector. These invariants are not yet the
full spectrum-level output, but they are the stable backbone needed for all
later projections.

\subsection*{2A. Consolidated dependency spine for v2.27.0}
For the next mathematical pass, the previous informal chain
\[
	ext{Ground Condition} \Longrightarrow 	ext{HC} \Longrightarrow 	ext{Closure} \Longrightarrow 	ext{Confinement}
\]
should be treated as a single dependency spine rather than as four loosely
adjacent slogans. The point of the present local pass is not to over-claim a
finished theorem package, but to normalize what the monolith already uses as
its operative order of dependence.

\paragraph{Lemma C (dependency spine lemma).}
Assume the triolic core, the standing orthogonality condition, and the
phase-shifted readout discipline. Then every admissible multi-representative
reading of one and the same body requires a compensating identification map;
this yields HC. Once HC is active, admissible state evolution is restricted to
those sectors that remain compatible with the common body; this yields closure.
Whenever closure is read at the level of the higher coupled container, the
resulting admissible class is a confinement class.

\paragraph{Proof sketch.}
The first implication is the content of Proposition~A: without compensation,
multiple shifted or dual descriptions of one body drift apart. The second is
Proposition~B: once HC acts as balancing constraint, the admissible dynamics can
no longer leave the common structural container. The last implication is a
change of reading level rather than a new mechanism: closure in the higher
coupled setting is exactly the confinement statement used downstream in the
septinion/seption language.

\paragraph{Operational reading.}
The practical value of this dependency spine is that later sections may appeal
to one normalized arrow instead of reintroducing the same chain piecemeal. In
particular, the later projector, transport, and Millennium-matrix sections
should be read as refinements of this spine, not as independent foundations.
This is also the correct mathematics-first reading for that archived public release:
first the spine is stabilized, then the surrounding operator language is
hardened.

\paragraph{Local closure target.}
For the purposes of the v2.26.x local sequence, this block counts as improved
only if all later uses of HC emergence, closure, and confinement can be traced
back to this normalized dependency spine without silent appeals to the
computational teaser architecture.

\subsection*{2B. Projector/closure anchoring of the dependency spine}
The next hardening step is to make explicit that the closure arrow in the
normalized spine is not merely verbal. It is read downstream through projector
language. In other words, the chain
\[
\text{Ground Condition} \Longrightarrow \text{HC} \Longrightarrow \text{Closure}
\]
should already be understood as carrying an operator-level consequence: once HC
has been introduced as the compensating identification map, admissible readings
must be projectively stabilized inside the common structural sector.

\paragraph{Lemma D (projector anchoring lemma).}
Assume the dependency spine of Lemma~C. Then every admissible downstream use of
closure may be rewritten as a projector-compatible statement: there exists a
sector projector $P_{\mathrm{adm}}$ such that the HC-balanced reading is
admissible only through the image of $P_{\mathrm{adm}}$, while noncompatible
continuations lie outside this image.

\paragraph{Proof sketch.}
If closure is the restriction to the common structural container, then a
projector description is simply the operator-level reading of that restriction.
The role of $P_{\mathrm{adm}}$ is not to invent a new mechanism, but to express
in normalized language that the HC-balanced state is retained only in the
closure-compatible sector. Thus projector logic is not parallel to closure
logic; it is the stabilized formal reading of closure once the argument moves
into the operator layer.

\paragraph{Downstream consequence.}
This anchoring is what allows the later projector/closure and transport sections
to be treated as genuine mathematical refinements of the dependency spine.
Ground Condition $\Rightarrow$ HC $\Rightarrow$ Closure no longer terminates in
a slogan, but feeds directly into the admissible projector sector on which the
later transport and Millennium-matrix mathematics are allowed to act.

\paragraph{Local closure target.}
For the v2.26.x sequence, this block counts as improved only if later projector
uses can be traced back to the dependency-spine-plus-projector reading, rather
than being invoked as an independent auxiliary formalism.

\subsection*{2C. Transport/Lorentz continuation of the dependency spine}
The next hardening step is to show that the normalized spine does not stop at
closure and projector admissibility. Once the admissible sector has been fixed,
the later transport layer must be read as a constrained continuation of the same
chain. In particular, Lorentz-transport language is not allowed to appear as an
independent dynamical decoration; it must preserve the HC-balanced,
projector-stabilized structural identity already enforced upstream.

\paragraph{Lemma E (transport continuation lemma).}
Assume the dependency spine of Lemma~C together with the projector anchoring of
Lemma~D. Then every admissible transport between downstream readout frames must
preserve the HC-balanced closure sector. Equivalently, Lorentz-type transport is
admissible only as a continuation on the image of the stabilized projector, not
as a map into an unconstrained external sector.

\paragraph{Proof sketch.}
By Lemma~D, admissibility is already restricted to the projector-compatible
closure image. A transport map acting on downstream representatives is therefore
allowed only if it respects that image. Otherwise the map would reintroduce a
break between closure logic and readout logic, contradicting the role of HC as
the compensating identification mechanism. Hence transport does not weaken the
spine; it inherits and preserves it.

\paragraph{Downstream consequence.}
This continuation is what permits the later Lorentz layer to be read in a
mathematics-first way: not as a second origin of structure, but as the
frame-translation discipline internal to the already stabilized closure sector.
It also prepares the next pass in which the Millennium matrix is treated as a
mathematical object acting only on admissibly transported states.

\paragraph{Local closure target.}
For the v2.26.x sequence, this block counts as improved only if later transport
claims can be referred back to the dependency spine through the combined chain
Ground Condition $\Rightarrow$ HC $\Rightarrow$ Closure $\Rightarrow$
Projector Stabilization $\Rightarrow$ Admissible Transport.


\subsection*{2D. Millennium-matrix realization of the dependency spine}
The next closure step is to prevent the Millennium matrix from floating as a
separate computational or architectural metaphor. If the mathematics-first
sequence is to hold, then the matrix must be read downstream from the already
stabilized chain of Ground Condition, HC emergence, closure, projector
stabilization, and admissible transport. In that sense the Millennium matrix is
not an independent origin of admissibility; it is the global matrix-level
realization of the closure-preserving compensator architecture already fixed by
the earlier lemmas.

\paragraph{Lemma F (Millennium-matrix realization lemma).}
Assume the dependency spine of Lemma~C, the projector anchoring of Lemma~D, and
the transport continuation of Lemma~E. Then the Millennium matrix may be
introduced as a mathematically admissible operator only on the HC-balanced,
closure-compatible, projector-stabilized, and transport-preserving state space.
Equivalently, its admissible role is not to generate a new unconstrained sector,
but to realize at matrix level the already normalized confinement-preserving
structural action.

\paragraph{Proof sketch.}
By Lemma~C the structural dependence runs from ground condition to HC and from
HC to closure. By Lemma~D this closure is stabilized as an admissible projector
image, and by Lemma~E every downstream transport must preserve that image.
Hence a matrix-level realization can be admitted only if it acts on that same
stabilized transported image. If it acted outside this chain, the Millennium
matrix would become an additional primitive rather than a realization of the
existing architecture, breaking the mathematics-first normalization. Therefore
its mathematically acceptable status is that of a global operator reading of the
already enforced confinement structure.

\paragraph{Proposition I (admissible matrix-sector invariance).}
Let $\mathcal M$ denote the Millennium matrix restricted to the HC-balanced,
closure-compatible, projector-stabilized, and transport-preserving sector
identified by Lemmas~C--F. Then $\mathcal M$ preserves that admissible sector;
in particular,
\[
\mathcal M\,P_{\mathrm{adm}} = P_{\mathrm{adm}}\,\mathcal M
\qquad	ext{on the stabilized transported image.}
\]
Hence the Millennium-matrix realization is mathematically acceptable only as a
sector-preserving global operator, not as a generator of a fresh unconstrained
layer.

\paragraph{Proof sketch.}
By Lemma~F the admissible introduction of $\mathcal M$ is already downstream
from closure, projector stabilization, and admissible transport. Proposition~H
states that the transport layer commutes with the admissibility projector on the
stabilized sector. Therefore the only mathematically consistent matrix-level
realization is one that acts within that same sector and preserves its image.
If $\mathcal M$ failed to preserve $P_{\mathrm{adm}}$, the matrix would reopen a
non-admissible branch and would no longer realize the normalized confinement
architecture. The commutation relation above is thus the matrix-level form of
admissible sector preservation.

\paragraph{Mathematical significance.}
This proposition upgrades the matrix discussion from a merely allowed operator
reading to a first invariance statement. The Millennium matrix is not only
introduced downstream from the dependency spine; it is constrained to preserve
the same admissible image that closure, projector stabilization, and transport
have already fixed. This gives the matrix chapter a clearer mathematical status
inside the normalization program.

\paragraph{Downstream consequence.}
This is the step that allows the later Millennium-matrix chapter to be retained
without overtaking the mathematics. The matrix can now be cited as a downstream
operator realization of the dependency spine, rather than as a free-standing
meta-object. It also prepares the next A-block pass, where trace and branch
mathematics can be read as acting under the matrix-level realization of the same
closure discipline.

\paragraph{Local closure target.}
For the v2.26.x sequence, this block counts as improved only if later uses of
the Millennium matrix can be referred back to the combined chain
Ground Condition $\Rightarrow$ HC $\Rightarrow$ Closure $\Rightarrow$
Projector Stabilization $\Rightarrow$ Admissible Transport $\Rightarrow$
Millennium-Matrix Realization.


\subsection*{2E. Trace/branch continuation of the dependency spine}
The final step in this first A-block pass is to prevent trace and branch
mathematics from re-opening the stabilized chain as if branching introduced a
new primitive source of admissibility. In the mathematics-first sequence, trace
formation and branch resolution must be read downstream from the already fixed
matrix-level realization. Thus branching is not a separate beginning of
structure; it is the controlled readout of differentiated admissible paths
inside the same confinement-preserving spine.

\paragraph{Lemma G (trace/branch continuation lemma).}
Assume the dependency spine of Lemma~C, the projector anchoring of Lemma~D, the
transport continuation of Lemma~E, and the Millennium-matrix realization of
Lemma~F. Then trace formation and branch resolution are mathematically
admissible only as downstream differentiations of the already HC-balanced,
closure-compatible, projector-stabilized, transport-preserving, and
matrix-realized state space. Equivalently, admissible trace/branch mathematics
does not create a fresh unconstrained dynamics; it resolves path structure
inside the same confinement-preserving architecture.

\paragraph{Proof sketch.}
By Lemma~F the Millennium matrix is admitted only as an operator realization of
the already stabilized dependency spine. Therefore any trace extracted from the
resulting evolution must be a trace of that same admissible operator action,
not a record of an independent unnormalized sector. Likewise, branch resolution
can be accepted only if it separates compatible downstream continuations of the
same confined structure. If branching were allowed to act as a new primitive,
then the chain from ground condition through HC, closure, projector
stabilization, transport, and matrix realization would no longer control the
later readout layer. Hence trace and branch mathematics remain admissible only
as terminal differentiations internal to the same spine.

\paragraph{Downstream consequence.}
This closes the first mathematics-first pass across the initial A-block group.
The trace/branch layer can now be cited without drifting into an autonomous
search or computational metaphor: it is anchored in the same structural chain
that begins at the ground condition and ends in admissible differentiated
readout. This is the point at which the first local freeze judgement for the
A-block group becomes meaningful.

\paragraph{Local closure target.}
For the v2.26.x sequence, this block counts as improved only if later trace and
branch statements can be referred back to the combined chain
Ground Condition $\Rightarrow$ HC $\Rightarrow$ Closure $\Rightarrow$
Projector Stabilization $\Rightarrow$ Admissible Transport $\Rightarrow$
Millennium-Matrix Realization $\Rightarrow$ Trace/Branch Continuation.


\paragraph{Proposition J (admissible trace/branch persistence).}
Let $\Gamma_{\mathrm{adm}}$ denote an admissible trace family generated inside the stabilized admissible sector, and let $B_{\mathrm{adm}}$ denote a branch selector that only separates continuations already induced by the same admissible matrix-sector action. Then admissible trace continuation and admissible branch separation preserve the same normalized admissible sector; in particular they do not generate a new independent admissibility source outside the projector/transport/matrix spine.
\[
B_{\mathrm{adm}}\,\Gamma_{\mathrm{adm}}\subseteq P_{\mathrm{adm}}\mathcal H_{\mathrm{conf}}.
\]

\paragraph{Proof sketch.}
By Proposition~H, admissible transport acts inside the stabilized projector image, and by Proposition~I the Millennium-matrix realization preserves that same admissible sector. Hence every admissible trace produced downstream is already a differentiated record of motion internal to the projector-stabilized, transport-compatible, matrix-preserving sector. If a branch operation is admissible at all, it cannot introduce a fresh primitive sector; it can only separate compatible continuations already contained in that same image. Therefore both trace continuation and admissible branching persist inside the normalized admissible sector.

\paragraph{Mathematical significance.}
This upgrades the final step of the first A-block from a narrative endpoint to a persistence statement: the trace/branch layer is no longer merely attached downstream, but mathematically constrained not to reopen an external admissibility source. In that sense the first A-block now carries a full preservation chain from the ground condition through HC, closure, projector stabilization, transport, matrix realization, and finally trace/branch continuation.

\paragraph{Local closure target.}
For the v2.26.x sequence, this block counts as strengthened only if later trace/branch arguments can explicitly cite Proposition~J as the persistence rule preventing downstream drift outside the admissible spine.

\subsection*{2F. Local freeze check for the projector/transport core}
The first A-block is now strong enough that its middle core --- projector,
closure, and transport --- can be evaluated on its own rather than only as a
passing stage inside the longer dependency spine. This local freeze check asks
whether that middle layer has reached a condition in which later work may treat
it as stable for the next release cycle without pretending that every possible
auxiliary lemma has already been unfolded to maximal density.

\paragraph{Local freeze status.}
At the present stage, the projector/closure plus transport/Lorentz core is best
read as \emph{locally freeze-ready for v2.27.0}, provided that no later pass
re-opens the admissible image as an unconstrained external sector. The closure
logic now fixes the admissible image; the projector logic stabilizes that
image; and the transport continuation restricts downstream frame changes to
that already fixed sector. In that sense, the central middle layer of the first
A-block is no longer merely planned architecture: it has acquired a coherent
mathematical dependency order.

\paragraph{Residual blocker.}
What still limits the block is not the absence of an admissibility spine, but
rather the remaining density gap between the local lemmas and a fully expanded
proof presentation. In other words, the structure is already in place, but the
companion could still benefit from additional explicit micro-lemmas later if a
future release wants a denser audit trail.

\paragraph{Release consequence.}
For the purposes of the mathematics-first roadmap, the projector/transport core
may therefore count as \emph{closed for now at medium proof density}. This is
enough to support the next local passes and to contribute positively to the
v2.27.0 gate, but not yet enough to justify claiming that every supporting
auxiliary statement has already reached final archival density.

\paragraph{Local closure target.}
The block counts as satisfactorily frozen for the next public release if later
passes preserve the chain
Ground Condition $\Rightarrow$ HC $\Rightarrow$ Closure $\Rightarrow$
Projector Stabilization $\Rightarrow$ Admissible Transport
without reintroducing a competing admissibility source outside that normalized
middle layer.



\subsection*{2G. Consolidated freeze and gate pass for the first A-block}
The first A-block has now been hardened not only by the dependency-spine lemmas,
but also by the three preservation propositions on transport, matrix action, and
trace/branch continuation. It is therefore possible to perform a consolidated
freeze and gate pass that asks whether the block, taken as a whole, has reached
a level suitable for counting positively toward the v2.27.0 release threshold.

\paragraph{Consolidated freeze status.}
Taken together, Lemma~C through Lemma~G and Proposition~H through Proposition~J
justify reading the first A-block as \emph{structurally freeze-ready for
v2.27.0 at medium proof density}. The block now carries a continuous admissible
spine from ground condition through HC, closure, projector stabilization,
transport, Millennium-matrix realization, and admissible trace/branch
persistence. No later local pass should need to reopen the block at the level
of its basic dependency order unless a genuine contradiction is found.

\paragraph{Gate consequence.}
For the global v2.27.0 gate, this means that the first A-block can now count as
a \emph{positive readiness contributor}. It does not by itself open the public
release, because the document-wide gate still depends on the remaining A-blocks,
reference stability, and hypothesis-control discipline. But it does remove the
first major blocker from the mathematics-first release logic.

\paragraph{Residual blocker.}
The remaining limitation is not structural incompleteness, but proof density and
audit granularity. Later archival passes may still unfold more micro-lemmas,
but this is no longer the same as saying that the block is mathematically open.
The correct current reading is: \emph{closed in structure, not maximally dense in
presentation}.

\paragraph{Local closure target.}
For the v2.26.x sequence, this pass counts as successful only if later global
readiness summaries can cite the first A-block as freeze-ready without again
collapsing its seven-step admissible spine into vague architectural language.

\subsection*{3. Higher embedding and confinement logic}
The local triolic organization is not the terminal description. Its load-bearing
role is to provide the minimal unit from which higher coupling becomes possible.
The octonionic step is useful because it already carries nontrivial coupled
imaginary directions, but within the present project logic it is not yet the
final confinement container. The septinion/seption reading is introduced
precisely to express that the entire compensated dynamics lives inside a more
rigid higher assembly in which uncontrolled sector break-out is blocked.

This point should be read carefully: the higher embedding is not decorative
algebraic escalation. It is the formal language by which the paper states that
closure is a global architectural fact rather than a local numerical accident.
In this sense, the passage
\[
\text{local triolic unit} \to \text{compensated coupling} \to \text{higher confinement}
\]
is one of the main structural arrows of the version.

\subsection*{4. First-layer master-equation logic}
The monolith already contains the structural core section and the ToE-style
master-equation section. v1.2.0 now fixes how these are to be read.
The first-layer master equation is not yet the maximal explicit final law of
the framework. It is the minimal compressed form that must simultaneously carry
geometry, triolic sector dynamics, compensator action, closure residue, and
observer-frame phase shift.

Accordingly, the five structural blocks used later in the paper should be
classified as follows.
\begin{enumerate}[leftmargin=2em]
  \item The geometric block is \emph{projectively privileged}: it is forced as a
  large-scale closure reading, but not as the primitive starting point.
  \item The triolic block is \emph{structurally forced}: without it there is no
  load-bearing internal organization.
  \item The HC block is \emph{structurally forced}: without it, no stable
  identification of dual/shifted representatives exists.
  \item The closure block is \emph{structurally forced}: without it, the theory
  would have no admissibility criterion.
  \item The observer-frame block is \emph{projectively forced}: once the
  measurement frame is shifted, a readout term is unavoidable.
\end{enumerate}

The upshot is that the master equation should always be read from left to right
as a compression chain,
\[
\text{structure} \to \text{compensation} \to \text{closure} \to \text{projection}.
\]
Any later physical naming --- Einstein-like, Dirac-like, or QFT-like --- comes
only after this structural order has been respected.

\paragraph{Proposition C (first-layer master form is forced).}
Once triolic organization, HC, closure, and frame shift are all imposed, the
dynamics can no longer be represented by a purely geometric or purely sectorial
equation. A combined master form is forced at first layer.

\paragraph{Proof sketch.}
Each omitted block destroys one already fixed architectural demand: omitting
triolicity erases the internal organization; omitting HC erases compensation;
omitting closure erases admissibility; omitting the observer term erases the
readout discipline. Hence only a combined structural form remains compatible.

\subsection*{5. Probe sector as second access path}
The probe sector is not a side remark but the second major access path of
v1.2.0. Two probes are sent toward one another, and each probe itself is
to be treated as a triolic structure. The third group stands orthogonal to both
probes and thereby defines a constrained search region rather than an arbitrary
ambient space. This orthogonal third group is crucial because it prevents the
probe architecture from degenerating into a mere two-body scan.

The role of the probe picture is twofold. First, it gives a structural way to
read interaction and search without abandoning the closure discipline. Second,
it provides an explicit place where the framework's search tools can be
attached. In the present version, the probe geometry is therefore a controlled
laboratory for structural read-off.

\paragraph{Proposition D (probe sector inherits closure).}
If each probe carries triolic organization and the third group remains
orthogonal to both probes, then the same HC/closure discipline that governs the
main sector also governs the admissible probe dynamics.

\paragraph{Proof sketch.}
The probe architecture does not introduce a new ontology. It duplicates the
same structural ingredients in a constrained interaction geometry. Since the
closure mechanism depends on these ingredients rather than on the label
``background'' or ``probe'', the probe sector inherits the same admissibility
logic.

\paragraph{Search space.}
The search space is therefore not an unrestricted continuum of candidates. It
is the region left after imposing triolic self-organization of each probe,
orthogonality of the third group, compensator compatibility, and closure.
Everything outside that region is structurally irrelevant for the probe readout.

\subsection*{6. Normalized tool apparatus}
A central goal of v1.2.0 is to stop treating rechen-tricks as scattered
intuition and instead store them as a reusable apparatus. The following classes
are fixed for this version.

\begin{description}[leftmargin=2.2em, style=nextline]
  \item[Closure-tools.] Tools that test, restore, or compare closure-compatible
  configurations. Their purpose is to decide admissibility and leakage.
  \item[Frame-shift-tools.] Tools that translate between the underlying
  structure and the observer's projected readout.
  \item[Projection-tools.] Tools that extract effective sector equations or
  recognizable physical forms from the common structural body.
  \item[Probe-tools.] Tools that organize the constrained search in the probe
  sector.
  \item[Equivalence-tools.] Tools that identify when two differently written
  states belong to the same structural class.
  \item[Read-off / trace-tools.] Tools that recover results from the structure
  without expanding the entire calculation tree every time.
\end{description}

Every individual tool should be documented under the same six labels:
\emph{name, purpose, prerequisite, short formal scheme, typical use, boundary
of validity}. This normalization rule is itself one of the deliverables of
v1.2.0.

\paragraph{Canonical examples in the present project.}
Kasisky detection is stored as a frequency/pattern detector rather than as a
cryptographic side story. Hash and rainbow methods are stored as routing and
candidate-compression devices. Markov organization is stored as a transition
controller over admissible search paths. Newton is stored as a convergence
accelerator once a structural branch has already been selected. GARCH, where
used at all, remains optional and secondary.

\subsection*{7. Projection discipline and physics reading}
The structural reading of matter, gauge, mass, and Higgs later in the monolith
is to be understood under a strict hierarchy.
\begin{enumerate}[leftmargin=2em]
  \item Matter as closure-stable bound sector is a structural reading.
  \item Gauge symmetry as invariance of closure-compatible transport is a
  projective reading.
  \item Mass as closure/frame residue is a projective reading with strong
  structural motivation.
  \item Higgs as coupling instance or mean coupling point is currently a
  hypothesis block and must remain explicitly marked as such.
\end{enumerate}

This hierarchy matters because it prevents later criticism from collapsing the
whole framework into one category. The paper does not claim that every named
physical object is already fully derived at the same level. It claims instead
that the structural origin is common and that different physical sectors appear
at different distances from that origin.

\subsection*{8. Consistency status of v1.2.0}
The present version should therefore be regarded as \emph{finished in working
form} whenever the following points are simultaneously true:
\begin{enumerate}[leftmargin=2em]
  \item the DOI, version stamp, and roadmap are fixed;
  \item the chain Ground Condition $\Rightarrow$ HC $\Rightarrow$ Closure is
  explicitly stated;
  \item the probe sector is presented as a genuine structural sector;
  \item the master equation is classified as a first-layer forced compression;
  \item the tool catalogue is normalized;
  \item open points are named rather than hidden.
\end{enumerate}

These conditions are now met in the present monolith together with the later
sections \emph{Structural Core and Closure Principles}, \emph{Toward a
Structural ToE Master Equation}, and \emph{Structural Reading of Matter,
Gauge, Mass, and Higgs}. What remains for v1.2.0 is not the recovery of a
missing architecture, but the next proof-theoretic deepening: more explicit
lemma chains, stronger sector-by-sector projections, and a larger appendix of
normalized tools and shortcuts.

